% GENERATED FILE. Edit canonical lessons, then run npm run generate:books.
% canonical-source-hash: fnv1a64:edf67e5b1d5a1407
% canonical-lessons: ZH-C277-youqu, ZH-C277-taoyan, ZH-C277-xinshang, ZH-C277-xiangshou, ZH-C277-reai, ZH-R277-review-272-277, ZH-R277-review-final-277

\chapter{Likes and Dislikes}
\label{ch:zh-likes-and-dislikes}
\hlchaptermodality{277}

\begin{chapteropening}
\textbf{By the end of this chapter:} \emph{I can say what is amusing, what I dislike, whom I admire, what I enjoy and what I love.}

Complete the last lesson of chapter 277: I can say what is amusing, what I dislike, whom I admire, what I enjoy and what I love.
\end{chapteropening}

\section[yǒuqù]{yǒuqù — fun; amusing}

\label{lesson:ZH-C277-youqu}



\begin{quote}
Before the new one: say the Mandarin for to retire, then the Mandarin for a big deal.
\end{quote}

\subsection*{You'll want to know: yǒuqù}
\emph{yǒuqù} — \textquotedblleft{}fun; amusing\textquotedblright{}.

\begin{quote}
\emph{zhège gùshi hěn yǒuqù} — This story is a lot of fun.
\end{quote}

Its characters come later; for now it is a word for the ear.

\begin{cousinweb}
\emph{yǒuqù} is built from \zh{有} 'to have' and \emph{qù}, both inherited from Middle Chinese, \emph{hjuwX} and \emph{tshjuH} (Guangyun, in Baxter's transcription), which glosses the second 'to tend towards'. How the pair came to mean amusing is not settled.
\end{cousinweb}

\begin{grammarlens}[title={What you've built}]
One more word for what I like.
\end{grammarlens}

\subsection*{Your turn}
\begin{itemize}
\raggedright
  \item \emph{Say it:} \emph{yǒuqù}
  \item \emph{Say it:} \emph{yǒuqù}, once more
  \item \emph{Say it:} \emph{zhège gùshi hěn yǒuqù} — This story is a lot of fun.
  \item \emph{Recall:} say the Mandarin for to retire, then the Mandarin for a big deal, then say \emph{yǒuqù} again
\end{itemize}

\begin{tcolorbox}[breakable,colback=teal!4,colframe=teal!35!black,title={Before you move on}]
What does yǒuqù mean? (\textquotedblleft{}Fun; amusing\textquotedblright{}.) Say it once more.
\end{tcolorbox}
\section[tǎoyàn]{tǎoyàn — to dislike; to hate}

\label{lesson:ZH-C277-taoyan}



\begin{quote}
Before the new one: say the Mandarin for a big deal, then the Mandarin for amusing.
\end{quote}

\subsection*{You'll want to know: tǎoyàn}
\emph{tǎoyàn} — \textquotedblleft{}to dislike; to hate\textquotedblright{}.

\begin{quote}
\emph{wǒ tǎoyàn xiàyǔ tiān} — I hate rainy days.
\end{quote}

Its characters come later; for now it is a word for the ear.

\begin{cousinweb}
\emph{tǎoyàn} is built from \emph{tǎo} and \emph{yàn}, both inherited from Middle Chinese, \emph{thawX} and \emph{'jiemH} (Guangyun, in Baxter's transcription), which glosses the first 'to govern' and 'to punish', and for the second quotes the Analects, on not tiring of fine food. How the pair came to mean disliking is not settled.
\end{cousinweb}

\begin{grammarlens}[title={What you've built}]
One more word for what I like.
\end{grammarlens}

\subsection*{Your turn}
\begin{itemize}
\raggedright
  \item \emph{Say it:} \emph{tǎoyàn}
  \item \emph{Say it:} \emph{tǎoyàn}, once more
  \item \emph{Say it:} \emph{wǒ tǎoyàn xiàyǔ tiān} — I hate rainy days.
  \item \emph{Recall:} say the Mandarin for a big deal, then the Mandarin for amusing, then say \emph{tǎoyàn} again
\end{itemize}

\begin{tcolorbox}[breakable,colback=teal!4,colframe=teal!35!black,title={Before you move on}]
What does tǎoyàn mean? (\textquotedblleft{}To dislike; to hate\textquotedblright{}.) Say it once more.
\end{tcolorbox}
\section[xīnshǎng]{xīnshǎng — to admire; to appreciate}

\label{lesson:ZH-C277-xinshang}



\begin{quote}
Before the new one: say the Mandarin for amusing, then the Mandarin for to dislike.
\end{quote}

\subsection*{You'll want to know: xīnshǎng}
\emph{xīnshǎng} — \textquotedblleft{}to admire; to appreciate\textquotedblright{}.

\begin{quote}
\emph{wǒ hěn xīnshǎng tā} — I really admire him.
\end{quote}

Its characters come later; for now it is a word for the ear.

\begin{cousinweb}
\emph{xīnshǎng} is built from \emph{xīn} 'glad' and \emph{shǎng} 'to reward', both inherited from Middle Chinese, \emph{xj+n} and \emph{syangX} (Guangyun, in Baxter's transcription), which glosses the first 'glad' and the second 'to bestow'.
\end{cousinweb}

\begin{grammarlens}[title={What you've built}]
One more word for what I like.
\end{grammarlens}

\subsection*{Your turn}
\begin{itemize}
\raggedright
  \item \emph{Say it:} \emph{xīnshǎng}
  \item \emph{Say it:} \emph{xīnshǎng}, once more
  \item \emph{Say it:} \emph{wǒ hěn xīnshǎng tā} — I really admire him.
  \item \emph{Recall:} say the Mandarin for amusing, then the Mandarin for to dislike, then say \emph{xīnshǎng} again
\end{itemize}

\begin{tcolorbox}[breakable,colback=teal!4,colframe=teal!35!black,title={Before you move on}]
What does xīnshǎng mean? (\textquotedblleft{}To admire; to appreciate\textquotedblright{}.) Say it once more.
\end{tcolorbox}
\section[xiǎngshòu]{xiǎngshòu — to enjoy}

\label{lesson:ZH-C277-xiangshou}



\begin{quote}
Before the new one: say the Mandarin for to dislike, then the Mandarin for to admire.
\end{quote}

\subsection*{You'll want to know: xiǎngshòu}
\emph{xiǎngshòu} — \textquotedblleft{}to enjoy\textquotedblright{}.

\begin{quote}
\emph{wǒmen zài xiǎngshòu jiàqī} — We're enjoying the holidays.
\end{quote}

Its characters come later; for now it is a word for the ear.

\begin{cousinweb}
\emph{xiǎngshòu} is built from \emph{xiǎng} and \emph{shòu}, both inherited from Middle Chinese, \emph{xjangX} and \emph{dzyuwX} (Guangyun, in Baxter's transcription), which lists the first under an older spelling and glosses it 'to offer' and 'to sacrifice', and glosses the second 'to take in'. How the first came to mean enjoying is not settled.
\end{cousinweb}

\begin{grammarlens}[title={What you've built}]
One more word for what I like.
\end{grammarlens}

\subsection*{Your turn}
\begin{itemize}
\raggedright
  \item \emph{Say it:} \emph{xiǎngshòu}
  \item \emph{Say it:} \emph{xiǎngshòu}, once more
  \item \emph{Say it:} \emph{wǒmen zài xiǎngshòu jiàqī} — We're enjoying the holidays.
  \item \emph{Recall:} say the Mandarin for to dislike, then the Mandarin for to admire, then say \emph{xiǎngshòu} again
\end{itemize}

\begin{tcolorbox}[breakable,colback=teal!4,colframe=teal!35!black,title={Before you move on}]
What does xiǎngshòu mean? (\textquotedblleft{}To enjoy\textquotedblright{}.) Say it once more.
\end{tcolorbox}
\section[rè'ài]{rè'ài — to love deeply}

\label{lesson:ZH-C277-reai}



\begin{quote}
Before the new one: say the Mandarin for to admire, then the Mandarin for to enjoy.
\end{quote}

\subsection*{You'll want to know: rè'ài}
\emph{rè'ài} — \textquotedblleft{}to love deeply\textquotedblright{}.

\begin{quote}
\emph{tā rè'ài yīnyuè} — He loves music.
\end{quote}

Its characters come later; for now it is a word for the ear.

\begin{cousinweb}
\emph{rè'ài} is built from \zh{热} 'hot' and \emph{ài} 'to love', both inherited from Middle Chinese, \emph{nyet} and \emph{'ojH} (Guangyun, in Baxter's transcription), which glosses the second 'to cherish'.
\end{cousinweb}

\begin{grammarlens}[title={What you've built}]
One more word for what I like.
\end{grammarlens}

\subsection*{Your turn}
\begin{itemize}
\raggedright
  \item \emph{Say it:} \emph{rè'ài}
  \item \emph{Say it:} \emph{rè'ài}, once more
  \item \emph{Say it:} \emph{tā rè'ài yīnyuè} — He loves music.
  \item \emph{Recall:} say the Mandarin for to admire, then the Mandarin for to enjoy, then say \emph{rè'ài} again
\end{itemize}

\begin{tcolorbox}[breakable,colback=teal!4,colframe=teal!35!black,title={Before you move on}]
What does rè'ài mean? (\textquotedblleft{}To love deeply\textquotedblright{}.) Say it once more.
\end{tcolorbox}
\section[(dialogue)]{Looking back: from Husband and Wife to Likes and Dislikes}

\label{lesson:ZH-R277-review-272-277}



\begin{quote}
Say the Mandarin for a wife and hubby.
\end{quote}

\subsection*{Your turn}
\begin{itemize}
\raggedright
  \item \emph{Say it:} a wife, hubby, grandpa on the mother's side, an uncle, an auntie
  \item \emph{Say it:} honest, hospitable, putting your heart into it, careless, careful with details
  \item \emph{Say it:} to ice-skate, to collect things, all the rage, handicrafts, a sports fan
  \item \emph{Say it:} to apologise, to forgive, to say yes to a request, to accept, as you like
  \item \emph{Say it:} to have a job interview, to quit a job, an opportunity, to retire, a big deal
  \item \emph{Say it:} amusing, to dislike, to admire, to enjoy, to love deeply
  \item \emph{Read it:} \zh{老} — each one aloud, then a word that holds it
\end{itemize}

\begin{tcolorbox}[breakable,colback=teal!4,colframe=teal!35!black,title={Before you move on}]
A wife? (\emph{qīzi}.) Hubby? (\textbf{\zh{老公}}, \emph{lǎogōng}.) Grandpa on the mother's side? (\textbf{\zh{外公}}, \emph{wàigōng}.) An uncle? (\emph{shūshu}.) An auntie? (\emph{āyí}.) Honest? (\emph{chéngshí}.) Hospitable? (\textbf{\zh{好客}}, \emph{hàokè}.) Putting your heart into it? (\textbf{\zh{用心}}, \emph{yòngxīn}.) Careless? (\emph{cūxīn}.) Careful with details? (\emph{xìxīn}.) To ice-skate? (\emph{huábīng}.) To collect things? (\emph{shōují}.) All the rage? (\textbf{\zh{热门}}, \emph{rèmén}.) Handicrafts? (\textbf{\zh{手工}}, \emph{shǒugōng}.) A sports fan? (\emph{qiúmí}.) To apologise? (\emph{dàoqiàn}.) To forgive? (\emph{yuánliàng}.) To say yes to a request? (\emph{dāying}.) To accept? (\emph{jiēshòu}.) As you like? (\emph{suíbiàn}.) To have a job interview? (\textbf{\zh{面试}}, \emph{miànshì}.) To quit a job? (\emph{cízhí}.) An opportunity? (\textbf{\zh{机会}}, \emph{jīhuì}.) To retire? (\emph{tuìxiū}.) A big deal? (\textbf{\zh{大事}}, \emph{dàshì}.) Amusing? (\emph{yǒuqù}.) To dislike? (\emph{tǎoyàn}.) To admire? (\emph{xīnshǎng}.) To enjoy? (\emph{xiǎngshòu}.) To love deeply? (\emph{rè'ài}.)
\end{tcolorbox}
\section[(dialogue)]{Likes and Dislikes — once more}

\label{lesson:ZH-R277-review-final-277}



\begin{quote}
Say the Mandarin for amusing and to dislike.
\end{quote}

\subsection*{Your turn}
\begin{itemize}
\raggedright
  \item \emph{Say it:} amusing, to dislike, to admire, to enjoy, to love deeply
\end{itemize}

\begin{tcolorbox}[breakable,colback=teal!4,colframe=teal!35!black,title={Before you move on}]
Amusing? (\emph{yǒuqù}.) To dislike? (\emph{tǎoyàn}.) To admire? (\emph{xīnshǎng}.) To enjoy? (\emph{xiǎngshòu}.) To love deeply? (\emph{rè'ài}.)
\end{tcolorbox}
